% Part: turing-machines
% Chapter: machines-computations
% Section: well-behaved-machines

\documentclass[../../../include/open-logic-section]{subfiles}

\begin{document}

\olfileid[tr]{tur}{mac}{dis}
\olsection{Disiplinli Makineler}

\begin{explain}
\olref[hal]{sec} bölümünde özel olarak belirlenmiş tek bir durma durumu~$h$
olan Turing makinelerini ele aldık; bu makinelerin durmaları hâlinde $h$
durumunda duracakları güvence altındadır. Bu bakımdan tek durma durumlu
makineler, genel olarak izin verdiğimiz Turing makinelerinden daha
``disiplinli''dir. Turing makinelerinin davranışına başka kısıtlamalar da
getirebiliriz. Örneğin Turing makinelerinin $0$. karedeki şerit sonu
işaretini silmesini ya da $0$. kareden sola gitmeye çalışmasını da
yasaklamadık. (Tanımımıza göre bu durumda kafa yalnızca $0$. karede kalır;
başka tanımlarda makine durur.) Benzer biçimde, bir Turing makinesinin
durması hâlinde $1$. karede durduğunu varsayabilmek bazen yararlıdır.
\end{explain}

\begin{defn}\ollabel{defn:disciplined}
Bir Turing makinesi~$M$ ancak ve ancak
\begin{enumerate}
    \item özel olarak belirlenmiş tek bir durma durumu~$h$ varsa,
    \item durması hâlinde $1$. kareyi tararken duruyorsa,
    \item $0$. karedeki $\TMendtape$ simgesini hiçbir zaman silmiyorsa ve
    \item $0$. kareden sola gitmeye hiçbir zaman çalışmıyorsa
\end{enumerate}
\emph{disiplinlidir}.
\end{defn}

\begin{explain}
Her Turing makinesinin, aynı davranışı gösteren ama özel bir durma durumu olan
bir makineye dönüştürülebileceğini daha önce tartıştık. Bunun için yalnızca
yeni bir $h$ durumu ve özgün $\delta$ fonksiyonunun tanımsız olduğu her
$\tuple{q,\sigma}$ çifti için
$\delta(q,\sigma)=\tuple{h,\sigma,N}$ yönergesi eklenir. Kanıtlamak zahmetli
olsa da her Turing makinesi~$M$'nin aynı girdilerde duran ve aynı çıktıyı
üreten disiplinli bir Turing makinesi~$M'$ biçimine getirilebileceği
doğrudur. Örneğin Turing makinesi durduğunda $1$. karede değilse kafanın şerit
sonu işaretini buluncaya kadar sola gitmesini, ardından bir kare sağa gidip
durmasını sağlayan yönergeler ekleyebiliriz. Öteki koşulların nasıl ele
alınabileceğini düşünmeyi size bırakıyoruz.
\end{explain}

\begin{ex}
    \olref{fig:adder-disc} içinde \olref[una]{ex:adder} içindeki toplama
    makinesini disiplinli bir makineye dönüştürüyoruz.
    \begin{figure}\[
\begin{tikzpicture}[->,>=stealth',shorten >=1pt,auto,node distance=2.8cm,
                    semithick]
  \tikzstyle{every state}=[fill=none,draw=black,text=black]

  \node[initial,state]         (A)              {$q_0$};
  \node[state]         (B) [right of=A] {$q_1$};
  \node[state]         (C) [below right of=B] {$q_2$};
  \node[state]         (D) [below left of=C] {$q_3$};
  \node[state]         (H) [left of=D] {$h$};

  \path (A) edge node {\TMtrans{\TMblank}{\TMstroke}{\TMstay}} (B)
           edge [loop below] node {\TMtrans{\TMstroke}{\TMstroke}{\TMright}} (B)
        (B) edge [loop below] node {\TMtrans{\TMstroke}{\TMstroke}{\TMright}} (B)
            edge node {\TMtrans{\TMblank}{\TMblank}{\TMleft}} (C)
        (C) edge node {\TMtrans{\TMstroke}{\TMblank}{\TMleft}} (D)
        (D) edge [loop above] node {\TMtrans{\TMstroke}{\TMstroke}{\TMleft}} (D)
        edge node {\TMtrans{\TMendtape}{\TMendtape}{\TMright}} (H);
\end{tikzpicture}
\]\caption{Disiplinli bir toplama makinesi}
\ollabel{fig:adder-disc}
\end{figure}
\end{ex}

\begin{prop}\ollabel{prop:disciplined} Her Turing makinesi~$M$ için öyle
disiplinli bir Turing makinesi~$M'$ vardır ki $M$, $O$ çıktısıyla duruyorsa
o da $O$ çıktısıyla durur; $M$ durmuyorsa o da durmaz. Özellikle, bir
Turing makinesince hesaplanabilen her $f\colon\Nat^n \to \Nat$ fonksiyonu
disiplinli bir Turing makinesince de hesaplanabilir.
\end{prop}

\begin{prob}\label{tur:mac:dis:prob:disc-succ}
$f(x)=x+1$ fonksiyonunu hesaplayan disiplinli bir makine verin.
\end{prob}


\begin{prob}\label{tur:mac:dis:prob:copier}
$\TMstroke^n$ girdisinde başlatıldığında
$\TMstroke^n \concat \TMblank \concat \TMstroke^n$ çıktısını üreten
disiplinli bir makine bulun.
\end{prob}

\end{document}
